\begin{claim}{B.25}{Proved}{Complete executed catalogues in types $(1,1)$ and $(0,4)$}
The extended mapping class groups of types $(1,1)$ and $(0,4)$ have,
respectively, $13$ and $30$ conjugacy classes of finite subgroups, including
the trivial subgroup.  The tables below give every class and the numbers of
full splittings, nontrivial partial splittings, maximal partial section
classes $\Phi$, and group orbits on $\Phi$.  The accompanying exact-arithmetic
files give every section, its image subgroup, its multiplier on every
element, its centralizer, all inclusion relations, the group stabilizers,
and the puncture stabilizers.  These are exhaustive catalogues for the two
specified surface types.
\end{claim}
\begin{proof}
Put
\[
 S=\begin{pmatrix}0&-1\\1&0\end{pmatrix},\qquad
 R=\begin{pmatrix}0&-1\\1&1\end{pmatrix},\qquad
 F=\begin{pmatrix}0&1\\1&0\end{pmatrix}.
\]
For the punctured torus use their linear actions on
$\R^2/\Z^2\setminus\{0\}$.  The homology map identifies its extended
mapping class group with $\mathrm{GL}_2(\Z)$: surjectivity is supplied by
linear maps, and a map acting trivially on homology can be isotoped to fix
two transverse basis curves, after which the complementary punctured disc
is handled by the Alexander isotopy.  This is the punctured version of
\cite[Theorem 2.5 and pp.~54--55]{bFarbMargalit2012}.

For the four-punctured sphere use the pillowcase
\[
 \bigl(\R^2/\Z^2\setminus (\tfrac12\Z^2/\Z^2)\bigr)/\{p\sim-p\}.
\]
Write $U(p)=p+(1/2,0)$ and $V(p)=p+(0,1/2)$.  The extended mapping
class group is
$V_4\rtimes\mathrm{PGL}_2(\Z)$ with these affine representatives.
The map to $\mathrm{PGL}_2(\Z)$ is obtained by lifting through the double
branched torus cover and acting on torus homology.  Its kernel consists of
the four half translations: a kernel element fixes all slope classes;
after a half translation it fixes the four punctures and two filling
curves, so the Alexander isotopy on their complementary discs kills it.
Linear maps split the homomorphism and act on half translations by
reduction modulo two.  This is the explicit morphism of
\cite[Proposition 2.7 and the extended statement on p.~57]{bFarbMargalit2012}.

Here are finite amalgams that make enumeration complete.  The matrices
$S,R$ generate $\mathrm{SL}_2(\Z)$ by the Euclidean algorithm, since
$S^{-1}R=\left(\begin{smallmatrix}1&1\\0&1\end{smallmatrix}\right)$.
Modulo $\{\pm I\}$ their images have orders two and three, with no
further relation: on the real projective line, $\bar S$ sends the negative
half-line to the positive half-line, and both nonidentity powers of
$\bar R$ send the positive half-line into the negative half-line.  The
free-product ping-pong argument gives
$\mathrm{PSL}_2(\Z)=C_2*C_3$.
Conjugation by $F$ fixes $\bar S$ and inverts $\bar R$.  Consequently
\[
 \mathrm{PGL}_2(\Z)=V_4*_{C_2}S_3,
 \qquad
 \mathrm{GL}_2(\Z)=\langle S,F\rangle
       *_{\langle-I,F\rangle}\langle R,F\rangle
       =D_8*_{V_4}D_{12},
\]
where $D_k$ has order $k$.  The second identity follows also by lifting
the first through their common central subgroup $\langle-I\rangle$;
the normal form for the amalgam leaves no additional kernel.
Adding the normal group of half translations gives the pillowcase amalgam
\[
 (V_4\rtimes V_4)*_{V_4\rtimes C_2}(V_4\rtimes S_3),
\]
whose vertex groups have orders $16$ and $24$, and whose edge group has
order $8$.  The two vertex groups are $D_8\times C_2$ and $S_4$.

A finite subgroup of an amalgam fixes a vertex in its Bass--Serre tree.
Indeed the convex hull of a finite orbit is a finite tree; its centre is
fixed, and the action preserves the two vertex types, so an edge centre
also gives fixed endpoints.  Two vertex subgroups are conjugate in the
amalgam precisely through the equivalence generated by conjugacy in
each vertex group and identification of edge subgroups.  To prove the
nontrivial direction, a conjugacy makes the subgroup fix both endpoints
of the intervening path; it fixes every edge on that path, which
decomposes the conjugacy into those successive vertex and edge steps.
Thus finite enumeration in the two vertices, followed by these exact
fusion steps, is a complete conjugacy calculation in the infinite group.
The program finds $10$ and $16$ actual subgroups in the torus vertices,
and $35$ and $30$ in the pillowcase vertices.  Fusion leaves exactly
$13$ and $30$ classes.  All subgroup lists are saturated under adjoining
an arbitrary vertex-group element, so the enumeration has no search bound
that could omit a subgroup.

We next justify the computation of sections.  Let
$Q=(\tfrac1{12}\Z^2)/\Z^2$.  Every isolated fixed point needed for
either model lies in $Q$.  For the torus it solves $(A-I)p=0$; for the
pillowcase it solves
\[
 (A-\epsilon I)p=-t/2\pmod{\Z^2},\qquad
 \epsilon\in\{1,-1\},\quad t\in(\Z/2)^2.
\]
For every positive matrix in the finite vertices, the inverse matrices
and the particular solutions of the nonsingular equations have
denominators dividing $12$.  The program certifies this by exact
fractions for every matrix and translation, rather than inferring it
from a grid count.  The singular case for a nonidentity positive element
is a nontrivial translation and contributes no solutions to that sign.

For a reflection, the same equations are affine circle equations.  The
greatest common divisor of the nonzero matrix entries is one or two, so
every component meets $Q$.  On each such circle the primitive integral
kernel direction, divided by $12$, joins consecutive sampled points.
In the pillowcase the plus and minus circles meet only at two-torsion
points, all of which are included in $Q$.  Quotienting the sampled graph
by $p\sim-p$ therefore gives a subdivision of the true fixed circles.
Deleting the prescribed punctures gives exactly their connected
components.  A component adjacent to a deleted point is a line; the
others are ovals.  All rotations and all punctures are retained as
vertices, so no incidence needed for restriction of sections is lost.

B.3 and B.4 identify these fixed components with section classes.
For a partial section with a rotation use its fixed point; for a
reflection use its fixed line or oval.  Inclusion of image subgroups and
containment of the fixed point in the fixed component determine exactly
when one section extends another.  Applying each affine element gives
the exact action on this finite section poset.  This computes $\Phi$,
its orbits and its stabilizers without replacing pointwise section data
by conjugacy of image subgroups in the quotient.  Centralizers are $1$
for isolated points and lines and $\Z$ for ovals, by B.3--B.4.

At a torus fixed point the local positive derivative is $A$.  At a
nonbranch pillowcase fixed point it is $\epsilon A$, with the unique
sign specified by $Ap+t/2=\epsilon p$.  For these matrices the traces
$-2,-1,0,1$ give orders $2,3,4,6$; the half-turn has multiplier $1/2$,
and otherwise the sign of $\det(e_1,Ae_1)$ selects $+1/n$ or $-1/n$.
Reflections have multiplier $1/2$.  At a removed pillowcase branch
point the quotient coordinate squares a torus coordinate, so the
positive rotation angle doubles.  The program records these puncture
data separately.  It checks the order and power identities for every
element of every section and checks that every section containing a
rotation has exactly one maximal overgroup.  The reflection classes
with several maximal extensions are retained explicitly.

Finally, C.5--C.6 show that the chosen finite actions compute the action
attached to each mapping class subgroup up to conjugacy, so these
fixed-component computations depend on the listed subgroup class, not
on an unrecorded realization choice.
\end{proof}

In the tables $f$ is the number of full splitting classes, $p$ the number
of nontrivial partial splitting classes with image subgroup retained,
$\phi=|\Phi|$, and $o=|\Phi/\Gamma|$.  Juxtaposition in the generator
column means composition of the matrices or affine maps defined above.
For $(0,4)$ the letters $S,R,F$ mean their pillowcase classes.
The complete output files are
\texttt{type\_1\_1\_finite\_subgroups.csv},
\texttt{type\_0\_4\_finite\_subgroups.csv}, and the corresponding
\texttt{partial\_sections.json} files, produced by
\texttt{catalogue\_low\_complexity.py}.

\begin{center}Type $(1,1)$\end{center}
\begin{longtable}{lllrrrrr}
\toprule
Class & Group & Generators & $|\Gamma^+|$ & $f$ & $p$ & $\phi$ & $o$\\
\midrule\endhead
T01 & $1$ & $1$ & 1 & 1 & 0 & 0 & 0\\
T02 & $C_2$ & $SF$ & 1 & 2 & 2 & 2 & 2\\
T03 & $C_2$ & $SSF$ & 1 & 1 & 1 & 1 & 1\\
T04 & $C_2$ & $SS$ & 2 & 3 & 3 & 3 & 3\\
T05 & $C_3$ & $RR$ & 3 & 2 & 2 & 2 & 2\\
T06 & $C_2^2$ & $SS,SF$ & 2 & 3 & 10 & 3 & 3\\
T07 & $C_2^2$ & $SS,SSF$ & 2 & 1 & 6 & 3 & 2\\
T08 & $C_4$ & $S$ & 4 & 1 & 4 & 3 & 2\\
T09 & $S_3$ & $RR,RRF$ & 3 & 2 & 7 & 2 & 2\\
T10 & $S_3$ & $RR,RF$ & 3 & 0 & 5 & 5 & 2\\
T11 & $C_6$ & $R$ & 6 & 0 & 5 & 5 & 2\\
T12 & $D_8$ & $S,SF$ & 4 & 1 & 15 & 3 & 2\\
T13 & $D_{12}$ & $R,RRF$ & 6 & 0 & 16 & 5 & 2\\
\bottomrule
\end{longtable}

\begin{center}Type $(0,4)$\end{center}
\begin{longtable}{lllrrrrr}
\toprule
Class & Group & Generators & $|\Gamma^+|$ & $f$ & $p$ & $\phi$ & $o$\\
\midrule\endhead
S01 & $1$ & $1$ & 1 & 1 & 0 & 0 & 0\\
S02 & $C_2$ & $SF$ & 1 & 4 & 4 & 4 & 4\\
S03 & $C_2$ & $SFV$ & 1 & 1 & 1 & 1 & 1\\
S04 & $C_2$ & $SFUV$ & 1 & 0 & 0 & 0 & 0\\
S05 & $C_2$ & $F$ & 1 & 2 & 2 & 2 & 2\\
S06 & $C_2$ & $V$ & 2 & 2 & 2 & 2 & 2\\
S07 & $C_2$ & $S$ & 2 & 0 & 0 & 0 & 0\\
S08 & $C_3$ & $RR$ & 3 & 1 & 1 & 1 & 1\\
S09 & $C_2^2$ & $V,SF$ & 2 & 2 & 9 & 4 & 3\\
S10 & $C_2^2$ & $V,SFU$ & 2 & 0 & 3 & 3 & 2\\
S11 & $C_2^2$ & $UV,SF$ & 2 & 0 & 6 & 6 & 3\\
S12 & $C_2^2$ & $UV,SFV$ & 2 & 2 & 6 & 2 & 2\\
S13 & $C_2^2$ & $UV,F$ & 2 & 2 & 8 & 2 & 2\\
S14 & $C_2^2$ & $SF,F$ & 2 & 0 & 6 & 6 & 3\\
S15 & $C_2^2$ & $SFUV,F$ & 2 & 0 & 2 & 2 & 1\\
S16 & $C_4$ & $FU$ & 2 & 0 & 2 & 2 & 1\\
S17 & $C_2^2$ & $V,U$ & 4 & 0 & 6 & 6 & 3\\
S18 & $C_2^2$ & $UV,S$ & 4 & 0 & 2 & 2 & 1\\
S19 & $C_4$ & $SU$ & 4 & 2 & 4 & 2 & 2\\
S20 & $S_3$ & $RR,FR$ & 3 & 1 & 8 & 4 & 2\\
S21 & $C_2^3$ & $V,U,SF$ & 4 & 0 & 18 & 6 & 3\\
S22 & $C_2^3$ & $UV,SF,F$ & 4 & 0 & 12 & 6 & 2\\
S23 & $C_4\times C_2$ & $FU,SU$ & 4 & 0 & 8 & 6 & 2\\
S24 & $D_8$ & $FU,V$ & 4 & 0 & 12 & 6 & 2\\
S25 & $D_8$ & $SU,SFV$ & 4 & 2 & 16 & 2 & 2\\
S26 & $D_8$ & $FU,SFV$ & 4 & 0 & 6 & 2 & 1\\
S27 & $D_8$ & $SU,V$ & 8 & 0 & 8 & 6 & 2\\
S28 & $A_4$ & $RR,RUR$ & 12 & 0 & 10 & 10 & 2\\
S29 & $D_8\times C_2$ & $FU,SU,V$ & 8 & 0 & 28 & 6 & 2\\
S30 & $S_4$ & $FUR,URF$ & 12 & 0 & 32 & 10 & 2\\
\bottomrule
\end{longtable}
